\begin{enumerate}

\item Identificar si la sección está abierta, cerrada unicelular o bicelular; dibujar su geometría y espesores.

\item Para una célula cerrada emplear Bredt, con q constante y tau=q/t; el espesor variable afecta al giro.

\item Comprobar que el par exterior coincide con el que transmiten los flujos, si la sección es cerrada.

\item Calcular el giro total con theta = M\_x L/(G J), o con la compatibilidad de las celdas.

\item Las magnitudes de la solución se han contrastado mediante cálculo analítico y el motor estructural de Studium cuando éste ofrece esa comprobación.

\end{enumerate}

\noindent\textbf{Operaciones y resultados intermedios.}

\begin{enumerate}

\item Área de línea media.

\[\fitmath{y_{1}=8100\approx 8100\,\mathrm{mm^{2}}}\]

\item Flujo cortante de Bredt.

\[\fitmath{y_{2}=\frac{900000}{2\cdot 90\cdot 90}\approx 55.55556\,\mathrm{N/mm}}\]

\item Constante de torsión cerrada.

\[\fitmath{y_{3}=\frac{4\cdot {\left(90\cdot 90\right)}^{2}}{\frac{90}{1.5}+\frac{90}{1.5}+\frac{90}{1.5}+\frac{90}{1.5}}\approx 1093500\,\mathrm{mm^{4}}}\]

\item Giro total entre extremos.

\[\fitmath{y_{4}=\frac{900000\cdot 800}{26955\cdot \frac{4\cdot {\left(90\cdot 90\right)}^{2}}{\frac{4\cdot 90}{1.5}}}\approx 0.02442724\,\mathrm{rad}}\]

\item Tensión tangencial máxima.

\[\fitmath{y_{5}=\frac{900000}{2\cdot 90\cdot 90\cdot 1.5}\approx 37.03704\,\mathrm{MPa}}\]

\end{enumerate}
